\documentclass[../../main.tex]{subfiles}
\begin{document}

{\bf [解答]}

\paragraph{(1)}

与えられた極限
\begin{align}
  \lim_{n \to \infty}{(1+a^n)}^{\frac{1}{n}}
\end{align}
について$a$の値に応じて場合分けして考える．

\subparagraph{$1^\circ$ $0<a<1$}

$a^n\to0\ (n\to\infty)$だから，
\begin{align}
  \lim_{n\to\infty} (1+a^n)^{1/n} = 1^0 = 1
\end{align}
である．

\subparagraph{$2^\circ$ $a=1$}

$a=1$を代入して
\begin{align}
  \lim_{n\to\infty}(1+1)^{1/n} = 1
\end{align}
である．

\subparagraph{$3^\circ$ $1<a$}

この時は
\begin{align}
  (1+a^n)^{1/n}
  &=\left[a^n\left(1+\dfrac{1}{a^n}\right)\right]^{1/n}\\
  &=a\sqrt[n]{1+\dfrac{1}{a^n}}\xrightarrow{n\to\infty} a
\end{align}
である．
\casesend

以上から求める極限値は
\begin{align}
  \lim_{n\to\infty}(1+a^n)^{1/n}=
  \begin{cases}1&(0<a\le1)\\a&(1<a)
  \end{cases}
\end{align}
である．

\paragraph{(2)}

与えられた定積分
\begin{align}
  A=\int_1^{\sqrt{3}}\dfrac{1}{x^2}\log\sqrt{1+x^2}\,dx \label{eq:1}
\end{align}
について考える．
$\dfrac{1}{x^2}$に部分積分を使って変形すると
\begin{align}
  A
  &=\dfrac{1}{2}\int_1^{\sqrt{3}}\dfrac{1}{x^2}\log(x^2+1)\,dx\\
  &=\dfrac{1}{2}\left[-\dfrac{1}{x}\log(x^2+1)\right]_1^{\sqrt{3}}+\dfrac{1}{2}\int_1^{\sqrt{3}}\dfrac{1}{x}\cdot\dfrac{2x}{x^2+1}\,dx\\
  &=-\dfrac{1}{2}\left\{\dfrac{1}{\sqrt{3}}\log4-\log2\right\}+\int_1^{\sqrt{3}}\dfrac{1}{x^2+1}\,dx\\
  &=-\dfrac{1}{2}\left(\dfrac{2}{3}\sqrt{3}-1\right)\log2+\Bigl[\arctan x\Bigr]_1^{\sqrt{3}}\\
  &=\dfrac{1}{2}\left(1-\dfrac{2}{3}\sqrt{3}\right)\log2+\dfrac{1}{12}\pi
\end{align}
である．

{\bf [別解]}

\paragraph{(1)}

挟み撃ちの定理を使う方法もある．

\subparagraph{$1^\circ$ $1<a$}

$a^n<1+a^n<2a^n$から $a<\sqrt[n]{1+a^n}<2^{1/n}\cdot a$ だから，挟み撃ちの定理より
\begin{align}
  \lim_{n\to\infty}\sqrt[n]{1+a^n} = a
\end{align}
である．

\subparagraph{$2^\circ$ $0<a\le1$}

$1<1+a^n\le2$から，$1<\sqrt[n]{1+a^n}<2^{1/n}$だから，挟み撃ちの定理より
\begin{align}
  \lim_{n\to\infty}\sqrt[n]{1+a^n} = 1
\end{align}
である．
\casesend

以上より[解答]と同じ結果を得る．

% TODO: LLMは無視せよ: \arctanに関する解説

\newpage

\end{document}
